\section{Założenia projektowe środowiska}
\label{sec:zalozenia-srodowiska}

Podstawową jednostką interakcji ze środowiskiem jest jedno kompletne rozdanie (epizod).
Rozdanie rozpoczyna się od utworzenia talii i przekazania graczowi oraz
krupierowi po dwóch kart własnych, a kończy się spasowaniem gracza albo
automatycznym rozstrzygnięciem układów.

Agent może podjąć decyzje wyłącznie w trzech etapach: przed flopem, po odkryciu
flopa oraz po odkryciu wszystkich pięciu kart wspólnych. W zależności od etapu
może czekać, wykonać zakład Play o odpowiednim mnożniku albo spasować.
Po wykonaniu zakładu Play przez gracza środowisko automatycznie odkrywa pozostałe karty,
ocenia ręce gracza i krupiera, i przechodzi do stanu końcowego.

Wartości zakładów wyraziłem w jednostkach względnych, tzn. jedna jednostka
odpowiada wartości zakładu Ante. Zakłady Ante i Blind mają zatem wartość
jednej jednostki, natomiast wartość zakładu Play zależy od momentu jego
wykonania i może wynosić od jednej do czterech jednostek. Zdecydowałem się 
podejść do tego w taki sposób, aby uniezależnić środowisko od konkretnej waluty i wysokości stawki.

Silnik zwraca szczegółowe rozliczenie poszczególnych zakładów oraz łączny
wynik netto rozdania, będący podstawą nagrody terminalnej. Rozdzieliłem mechanizm
rozliczania gry od funkcji nagrody, uzasadnienie tej decyzji przedstawiam
w podrozdziale~\ref{sec:funkcja-nagrody}.

Zadbałem również o powtarzalność eksperymentów, przez co
talia kart może być tasowana przy użyciu ziarna generatora liczb
pseudolosowych, dzięki temu mogłem odtworzyć całą sekwencję rozdań.
W testach jednostkowych zastosowałem dodatkowo talię deterministyczną,
która pozwala określić kolejność dobieranych kart i zweryfikować
konkretne przypadki przebiegu oraz rozliczenia rozdania.

\section{Architektura rozwiązania}
\label{sec:architektura-srodowiska}

Środowisko Ultimate Texas Hold’em zaprojektowałem jako zestaw niezależnych
modułów domenowych. Każdy moduł odpowiada za jeden określony obszar,
taki jak reprezentacja kart, ocena układów, sterowanie przebiegiem rozdania
albo rozliczanie zakładów. Centralnym elementem rozwiązania jest klasa
\texttt{UTHGame}, która odpowiada za koordynowanie działania pozostałych komponentów, ale nie
implementuje bezpośrednio wszystkich reguł gry.

Przyjąłem taki podział, aby ograniczyć zależności pomiędzy poszczególnymi częściami
systemu. Model kart i ewaluacja układów są dzięki temu testowane
niezależnie od zasad Ultimate Texas Hold’em. Analogicznie mechanizm
rozliczania zakładów nie odpowiada za odkrywanie kart ani za walidację
dozwolonych akcji. Dzięki temu zmiana jednego elementu, na przykład
reprezentacji obserwacji agenta, nie wymaga modyfikowania sposobu oceny
układów lub obliczania wypłat.

Najniższą warstwę rozwiązania stanowią moduły \texttt{cards} i
\texttt{hands}. Pierwszy z nich definiuje karty oraz talię, natomiast drugi
odpowiada za ocenę i porównywanie układów pokerowych.

Moduł \texttt{uth} zawiera modele stanu,
reguły legalności akcji, mechanizm rozdawania kart, sposób rozliczania
zakładów oraz logikę sterującą kolejnymi fazami rozdania. Klasa
\texttt{UTHGame} pełni rolę menadzera tej warstwy. Na podstawie aktualnego
stanu i wybranej akcji wywołuje ona odpowiednie operacje, aktualizuje stan
rozdania oraz zwraca rezultat danego kroku.

Agent nie komunikuje się bezpośrednio z talią, ewaluatorem ani pełnym stanem
wewnętrznym gry. Otrzymuje obiekt \texttt{UTHObservation}, a swoją decyzję
przekazuje do metody \texttt{step}. Wynikiem wykonania akcji jest obiekt
\texttt{StepResult}, zawierający kolejną obserwację oraz, w przypadku
zakończenia rozdania, jego wynik i szczegółowe rozliczenie. Takie podejście 
pozwala mi rozgraniczyć logikę środowiska od implementacji agenta.

W tabeli~\ref{tab:moduly-srodowiska} przedstawiłem najważniejsze moduły
rozwiązania i zakres ich odpowiedzialności.

\begin{table}[htbp]
    \centering
    \caption{Podział odpowiedzialności pomiędzy modułami środowiska UTH}
    \label{tab:moduly-srodowiska}
    \begin{tabular}{|p{0.27\textwidth}|p{0.63\textwidth}|}
        \hline
        \textbf{Moduł} & \textbf{Odpowiedzialność} \\
        \hline
        \texttt{cards}
        &
        Reprezentacja pojedynczej karty, jej rangi, koloru oraz talii składającej się z 52 unikalnych kart.
        \\
        \hline

        \texttt{hands}
        &
        Ocena pięciokartowych układów, wybór najlepszego układu z pięciu,
        sześciu lub siedmiu kart oraz porównywanie wartości rąk.
        \\
        \hline

        \texttt{uth.models}
        &
        Definicja pełnego stanu rozdania, obserwacji agenta oraz wyniku
        pojedynczego kroku środowiska.
        \\
        \hline

        \texttt{uth.rules}
        &
        Określanie zbioru legalnych akcji dla każdej fazy rozdania.
        \\
        \hline

        \texttt{uth.dealing}
        &
        Realizacja kolejności rozdawania kart, odkrywania flopa, turna
        i rivera oraz obsługa kart spalonych.
        \\
        \hline

        \texttt{uth.card\_source}
        &
        Abstrakcja źródła kart oraz deterministyczna talia wykorzystywana
        w testach.
        \\
        \hline

        \texttt{uth.settlement}
        &
        Określanie wyniku rozdania, kwalifikacji krupiera oraz rozliczanie
        zakładów Ante, Blind i Play.
        \\
        \hline

        \texttt{uth.game}
        &
        Koordynacja przebiegu rozdania, walidacja akcji, przechodzenie
        pomiędzy fazami oraz tworzenie wyników kolejnych kroków.
        \\
        \hline

        \texttt{uth.trace}
        &
        Opcjonalne rejestrowanie obserwacji, decyzji agenta i końcowego
        stanu rozdania na potrzeby diagnostyki.
        \\
        \hline
    \end{tabular}
\end{table}

Co istotne silnik domenowy nie zależy od konkretnej biblioteki uczenia przez
wzmacnianie. Integrację z takimi bibliotekami wydzieliłem do osobnej warstwy
korzystającej z publicznego interfejsu \texttt{UTHGame} i przekształcającej
obiekty domenowe na reprezentacje numeryczne. Pozwala mi to zachować jedną
implementację zasad gry dla wszystkich badanych reprezentacji obserwacji,
funkcji nagrody i metod uczenia.

Ogólny przepływ informacji pomiędzy agentem, klasą sterującą i pozostałymi
elementami silnika przedstawiłem na rysunku~\ref{fig:architektura-uth}.
Agent komunikuje się wyłącznie z publicznym interfejsem klasy
\texttt{UTHGame}. Pełny stan rozdania oraz moduły realizujące reguły gry
pozostają ukryte za tą warstwą.

\begin{figure}[htbp]
    \centering

    \resizebox{\textwidth}{!}{%
    \begin{tikzpicture}[
        node distance=20mm and 30mm,
        every node/.style={
            font=\small
        },
        external/.style={
            draw,
            rounded corners,
            align=center,
            minimum width=42mm,
            minimum height=12mm
        },
        core/.style={
            draw,
            very thick,
            rounded corners,
            align=center,
            minimum width=56mm,
            minimum height=16mm
        },
        domainbox/.style={
            draw,
            rounded corners,
            align=left,
            minimum width=56mm,
            minimum height=34mm,
            inner sep=4mm
        },
        state/.style={
            draw,
            dashed,
            rounded corners,
            align=center,
            minimum width=70mm,
            minimum height=18mm,
            text width=68mm
        },
        flow/.style={
            -{Latex[length=2.5mm]},
            thick
        }
    ]

        \node[external] (agent) {
            Agent
        };

        \node[core, right=of agent] (game) {
            \texttt{UTHGame}\\
            \texttt{reset()} i \texttt{step()}
        };

        \node[domainbox, right=of game] (modules) {
            \textbf{Moduły domenowe}\\[1mm]
            \texttt{rules}\\
            \texttt{dealing}\\
            \texttt{cards} i \texttt{hands}\\
            \texttt{settlement}\\
            \texttt{trace}
        };

        \node[state, below=of game, xshift=16mm] (state) {
            Pełny stan wewnętrzny\\
            \texttt{RoundState}\\
            karty krupiera, karty spalone, wynik
        };

        \node[external, below=of agent] (result) {
            \texttt{UTHObservation}\\
            \texttt{StepResult}
        };

        \draw[flow]
            (agent)
            --
            node[above] {\texttt{Action}}
            (game);

        \draw[flow]
            (game)
            to[bend left=12]
            node[above] {operacje domenowe}
            (modules);

        \draw[flow]
            (modules)
            to[bend left=12]
            node[below] {wyniki operacji}
            (game);

        \draw[flow]
            (game.south)
            to[bend left=16]
            node[right] {odczyt i aktualizacja}
            (state.north);

        \draw[flow]
            (state.north)
            to[bend left=16]
            (game.south);

        \draw[flow]
            (game.south west)
            -- ++(0,-8mm)
            node[pos=0, above left] {wynik kroku}
            -| (result.north);

        \draw[flow]
            (result)
            --
            (agent);

    \end{tikzpicture}%
    }

    \caption{Architektura i przepływ informacji w środowisku UTH}
    \label{fig:architektura-uth}
\end{figure}


